\section{Detector model and simulation chain (P2)}\label{sec:detector}

The P2 study asks what survives of the truth-level quirk patterns in a realistic tracker at pile-up 200. This section
describes the detector model, the pile-up overlay, the digitisation and the standard tracking that removes helical tracks.
Everything is built from public information: the CMS Phase-2 Tracker~\cite{CMS:2017lum}, MIP Timing
Detector~\cite{Butler:2019rpu} and Level-1 Trigger~\cite{Zabi:2020gjd} Technical Design Reports of the CMS
detector~\cite{CMS:2008xjf} and the public tkLayout
geometry. The design is in \rep{acts/stage4/plan.md} (reviewed in nine rounds), the frozen definitions in
\rep{acts/stage4/WP_RULE.md}, and the validation in the G1, G1b and G3 reports under \rep{acts/stage4/}.

\subsection{Chain overview}\label{sec:detector:chain}
The Stage-3 truth hits (configuration with field and nominal material, smeared luminous-region vertex, all surfaces,
$|\eta_{\rm pair}|<4$ propagated) are injected, one ACTS event per physics event, into a pile-up event simulated on the
same surfaces with ACTS~\cite{Ai:2021ghi} (pyacts 48.0.0) and its fast simulation Fatras. Quirk and pile-up hits are then
digitised together, a standard combinatorial Kalman filter (CKF) finds helical tracks, and a dedicated finder searches
the remaining hits for planes through the primary vertex (Section~\ref{sec:strategy}). Pile-up events are not stored:
each is regenerated bit for bit from the job seed and the event index.

\subsection{Geometry and material}\label{sec:detector:geo}
Each of the 83 physical Stage-3 surfaces except the endcap timing layer (ETL) becomes a sensitive ACTS surface at its nominal position in the D111 layout:
cylinders for the barrel (four IT pixel layers TBPX, three PS layers TBPS, three 2S layers TB2S, the barrel timing layer
BTL) and discs for the endcaps. The Stage-3 material per layer (Section~\ref{sec:p1:prop}) is attached as surface
material; stacked layers split it over their two sensors. Every cylinder and disc is divided into 32 azimuthal sectors of
the same geometric surface, which leaves the hit positions unchanged and brings the CKF from 23--60\,s to 1.6--3.4\,s per event. The
ETL could not be included in the ACTS geometry (a volume-stacking limitation); no finder variable uses it. A muon gun gives the full
17-hit central pattern with every hit on its surface to $<0.1\um$, and the Fatras scattering angle behind the first pixel
layer agrees with the Highland formula~\cite{Highland:1975pq} (ratio 0.93 with the PDG definition). Two Fatras pitfalls
were found and fixed: the default minimum momentum of $0.5\GeV$ drops every low-momentum looper (set to $30\MeV$), and
the JSON material density must be given in ACTS-native units.

\subsection{Pile-up overlay and density normalisation}\label{sec:detector:pu}
Minimum-bias interactions are generated with Pythia~8.312~\cite{Bierlich:2022pfr} (CP5~\cite{CMS:2019csb},
$\sqrt s=14\TeV$, $\sigma_{\rm inel}=78.9$\,mb), with a Poisson number of mean 200 per event and vertices drawn from the
luminous region ($\sigma_z=45\mm$, $\sigma_t=190\ps$~\cite{Butler:2019rpu}; $\sigma_{xy}=5\um$, \EST{}).
$K^0_S$, $\Lambda^0$ and other hyperons are decayed by Pythia; charged particles and photons with $p_T>50\MeV$ and $|\eta|<4$
are passed to Fatras with multiple scattering, ionisation, bremsstrahlung and photon conversions. Fatras has no hadronic
interactions and does no decays in flight, and out-of-time pile-up is not simulated.

The resulting in-time density, $1.33\times10^5$ hits per event, is \emph{below} the Tracker TDR occupancies (which
include out-of-time pile-up and secondaries~\cite[Sec.~6.2.1, Fig.~6.3]{CMS:2017lum}) by a factor 2.3--3.7 in the IT
(the IT cluster size is not given in the TDR, \EST{}), 2.5--4.3 in the PS and 4.0--7.3 in the 2S layers, rising with
radius. Since a fake-rate estimate at too low a density would be optimistic, every pile-up event is \textbf{topped up}
per layer to the TDR occupancy (user decision, Appendix~\ref{app:limits}): $({\rm ratio}-1)$ times the event's own Fatras crossings are added
on each layer, with $z$ resampled from the layer's hits and uniform $\phi$; on stacked layers they are added as in/out
pairs, the conservative choice for the per-crossing counts. Each added crossing counts as its own interaction in the
pile-up extrapolation (Section~\ref{sec:bkg:pu}). The lower and upper ends of the ratio ranges are a systematic.

\subsection{Digitisation, merging, timing and efficiency}\label{sec:detector:digi}
Hits are digitised by Gaussian smearing with binary resolutions (pitch$/\sqrt{12}$, as the TDR confirms for the
OT~\cite[Sec.~12.2]{CMS:2017lum}), with ACTS cluster merging switched off; two measurements on the same sensor are merged
at analysis level when they lie within one pitch in both local directions (user decision, Appendix~\ref{app:limits}), with the mean
position, an error of $\max(\sigma,|\Delta|/\sqrt{12})$ and summed energy. This sets the hit-level merging of the two
quirks at high $\Lam$. Table~\ref{tab:digi} lists the resolutions. Quirk hits arriving later than $\Delta t_{\rm
acc}=25\ns$ after a prompt particle would (one bunch crossing, following Ref.~\cite[Sec.~IV]{Knapen:2017kly}; variation
$12.5\ns$) are dropped, and a flat 99\% hit efficiency is applied (\UNV{}; variation 97\%).

\begin{table}[htbp]
\centering
\caption{Digitisation model (binary resolution, pitch$/\sqrt{12}$; \rep{plan.md} \S1.5). Pitches from the Tracker
TDR~\cite{CMS:2017lum} and the public D111 module table; BTL from the MTD TDR~\cite{Butler:2019rpu}.}
\label{tab:digi}
\small
\begin{tabularx}{\linewidth}{lccc>{\raggedright\arraybackslash}X}
\toprule
module & pitch & $\sigma_{r\phi}$ & $\sigma_{\rm long}$ & variation\\
\midrule
IT pixel & $25\times100\,\mu$m$^2$ & $7.2\um$ & $28.9\um$ & charge-sharing 3--6 / 11--24$\um$\\
PS (strip/macro-pixel) & $100\um\times1.467\mm$ & $28.9\um$ & $423\um$ & $\pm20\%$\\
2S (strips) & $90\um\times50.27\mm$ & $26.0\um$ & $14.5\mm$ & uniform smearer\\
BTL (timing) & bars $57\mm\times3.12\mm$ & $16\mm$ (\EST{}) & $0.9\mm$ & $\sigma_t=35\ps$; 60\,ps end of life\\
\bottomrule
\end{tabularx}
\end{table}

\subsection{Standard tracking and hit removal}\label{sec:detector:ckf}
Quirk tracks are not helices, so the CKF is used only to remove the hits of helical tracks before the quirk finder runs:
seeds from triplets of IT space points with $p_T>0.9\GeV$, $|z_0|<200\mm$, tracks with $\ge7$ measurements, $\le2$ holes
and $|d_0|<4\mm$, following the ``unassociated hits'' logic of Ref.~\cite[Sec.~III.D]{Knapen:2017kly}. In the production chain 13.9\% of all
measurements (19.5\% of the barrel hits, $\approx1540$ tracks per event) are removed; low-$p_T$ pile-up hits stay and form the combinatorial background. The fraction $f_{\rm CKF}$
of a candidate's band hits that sit on CKF tracks is a selection variable (Section~\ref{sec:strategy:finder}); on the probe
samples no quirk measurement was taken by a CKF track (0 of 208).

\subsection{Primary vertex, MET and trigger model}\label{sec:detector:pvmet}
The primary vertex is the true hard-scatter vertex smeared by $(\sigma_{xy},\sigma_z)=(10,15)\um$ with a 93\%
identification efficiency at pile-up 200 (graphical readings of the Tracker TDR~\cite{CMS:2017lum}, \UNV{}). The missing
transverse momentum is the truth $p_T(Q\bar Q)$ (the quirks deposit $\lesssim5\GeV$ in the calorimeters, carried as a
one-sided $\lesssim2\%$ bias) smeared by $\sigma_\perp=25\GeV$ per component (15 and $35\GeV$ variations; the PUPPI MET
resolution at pile-up 200 is \UNV{}). The trigger and preselection model is part of the selection and is given in
Section~\ref{sec:strategy:trig}.

\subsection{Samples}\label{sec:detector:samples}
Table~\ref{tab:samples} lists the detector-level samples. Production ran on one laptop (Apple M3, 8 cores) in three
campaigns of about one day each (C1a, C1b, C2B) plus a full-chain re-run of nine signal points (A2FC): 2295 jobs, no job
errors, $\approx382$ CPU-hours. About 0.14\% of pile-up events hit an ACTS navigation assertion at the $\phi$-sector
boundaries; they are replaced deterministically by a re-seeded event and recorded (three G3 pilot jobs that failed this way
before the fix were left out).

\begin{table}[htbp]
\centering
\small
\caption{Detector-level samples (\rep{acts/stage4/C1_LAUNCH.md}, \rep{G1B_REPORT.md}, \rep{A2FC_REPORT.md}).
$\rho$ is the pile-up density in units of a normalised PU-200 event (Section~\ref{sec:bkg:pu}).}
\label{tab:samples}
\begin{tabularx}{\linewidth}{lX}
\toprule
sample & content and purpose\\
\midrule
A1 & signal without pile-up: finder-efficiency calibration and the window/tolerance variations\\
A2 & signal + PU 200 on five (species, mass) points ($E$ 200, 500, 1000, 1500; $\Etil$ 500) $\times$ 8 $\Lam$ values, run row-major
     (the same pile-up event for every $\Lam$ of a row), $\Lam=0.45,0.6,0.85,1,1.5,2,3,5\keV$, 180--200 template events
     per point\\
A5 & five (species, mass) points at $\Lam=1\keV$ through the full ACTS chain with dumps (closure of the fast signal path)\\
A2FC & full-chain re-run of the nine A2 points where the fast signal path was not validated (pre-registered revert rule)\\
B1 & pile-up only, $\rho=1$, four MET roads per event: 12\,000 events, 48\,000 trials\\
B2 & pile-up only, $\rho=2$ (superposition of two normalised events): 3000 events, 6000 trials\\
dens & $\rho=1$ re-runs of a B1 subset with the top-up at its lower and upper ratio (density systematic)\\
SM & Pythia~8 $Z(\nu\nu)$+jet, $W(\ell\nu)$+jet (incl.\ $\tau$ split by decay class), $t\bar t$, with PU 200 and truth MET$>200\GeV$;
     2000 events per sample\\
C$'_0$, C$'_{\rm PU}$ & guns and toys for single-interaction topologies (photon conversions, $V^0$, vertex pairs, loopers,
     nuclear interactions, decays in flight): $3\times10^4$--$10^5$ per gun or toy without pile-up, 500 per topology (8 topologies) with
     pile-up\\
\bottomrule
\end{tabularx}
\end{table}
